	\documentclass[14pt]{extreport}
\usepackage{extsizes}
	\usepackage[frenchb]{babel}
	\usepackage[utf8]{inputenc}  
	\usepackage[T1]{fontenc}
	\usepackage{amssymb}
	\usepackage[mathscr]{euscript}
	\usepackage{stmaryrd}
	\usepackage{amsmath}
	\usepackage{tikz}
\usepackage{eurosym}
	\usepackage[all,cmtip]{xy}
	\usepackage{amsthm}
	\usepackage{varioref}
	\usetikzlibrary{patterns}
	\usepackage{float}
	\usepackage[ margin=.75in]{geometry}
	\geometry{a4paper}
	\usepackage{lmodern}
	\usepackage{hyperref}
	\usepackage{array}
	\usepackage{easytable}
	 \usepackage{fancyhdr}\usepackage{longtable}

	\pagestyle{fancy}
	\theoremstyle{plain}
	\fancyfoot[C]{\empty} 
	\fancyhead[L]{Contrôle}
	\fancyhead[R]{8 décembre 2025}
	
	
	\title{Contrôle chapitre 4}
	\date{}
	\begin{document}

\begin{center}\textbf{Soignez votre présentation et votre rédaction.}\end{center}

\subsection*{Exercice 0 (1 point)} 

Donner une valeur approchée de $\pi$ au centième. Que représente ce nombre ?
\subsection*{Exercice 1 (6 points)}  % 4 points

Faire une figure à main levée, et donner le périmètre des figures suivantes : \begin{enumerate}
\item Un losange dont un côté mesure $2$ cm. 
\item Un rectangle de longueur $3,2$ dm et de largeur $37$ cm. 
\item Un cercle de rayon $4$ cm (donner la valeur exacte, puis une valeur approchée au mm).
\item Le triangle ci-dessous. 
\begin{figure}[H]\center
\begin{tikzpicture}[scale = 3]
\draw (-0.5, 0) --  (3.5, 0) -- (0, 1.2) -- (-0.5, 0);
\draw[dashed] (0,0) -- (0,1.2); 
\draw [<->] (-0.5, -.1) -- (0, -.1); 
\draw [<->] (3.5, -.1) -- (0, -.1); 
\draw (-.25, -.2) node {$1,5$ cm};
\draw (1.65, -.2) node {$10,5$ cm};
%\draw [<->] (0.05, 0) -- (0.05, 1.2); 
\draw (0.29, .6) node {$3,6$ cm};
\draw (1.6, .8) node {$11,1$ cm};
\draw (-0.6, .6) node {$3,9$ cm};
\fill (0.1,0) -- (0.1, 0.1) -- (0, 0.1) -- (0,0);
\end{tikzpicture}
\end{figure}
\end{enumerate}

\subsection*{Exercice 2 (4 points)}

Calculer le périmètre des deux figures suivantes.
\[\includegraphics[width = 18 cm]{Exo1}\]



\newpage

\subsection*{Exercice 3 (4 points)}

On considère les deux figures suivantes :
\[\includegraphics[width = 18 cm]{Exo2}\]

Laquelle a le plus grand périmètre ? Détailler votre raisonnement.

\subsection*{Exercice 4 (5 points)}
On considère les deux figures suivantes : 

\begin{figure}[H]
\center 
\begin{tikzpicture}[scale =1.3]
\draw (-2, 0) -- (2, 0) arc (0: 180 :2);
\draw (-2, 0) -- (2, 0) arc (0: 180 :2);
\draw[<->] (-2, -.2) -- (2, -.2); 
\draw (0, -.5) node {$4\ cm$};
\draw[white](0,-1)--(3,-1);
\end{tikzpicture}
\begin{tikzpicture}[scale =1.3]
\draw[dashed](-2, 0) -- (2, 0);
\draw(2,0) arc (0: 180 :2) arc (180:360: 1) arc(180:0:1);
\draw(2,0) arc (0: 180 :2) arc (180:360: 1) arc(180:0:1);
\draw[<->] (2, -.2) -- (0, -.2); 
\draw (1, -.5) node {$2\ cm$};
\end{tikzpicture}
\end{figure} 
\begin{enumerate}
\item Calculer le périmètre des deux figures. Donnez d'abord une valeur exacte faisant apparaître $\pi$, puis une valeur approchée au millimètre.

\item Laquelle des figures a le plus grand périmètre ? 
\end{enumerate} 
 

\end{document}